\section{Evidências disponíveis e plano de testes}
\label{sec:avaliacao}

\subsection{Resultados computacionais do modelo}

As métricas da Tabela~\ref{tab:metricas} são provenientes da avaliação
registrada para o modelo compacto, nos lotes de teste descritos na
Seção~\ref{sec:modelo}. Elas avaliam correspondência com rótulos do
dataset local, não detecção experimental de fungos no mini-silo.

\begin{table}[htbp]
\centering
\caption{Avaliação do modelo no conjunto de teste}
\label{tab:metricas}
\begin{tabularx}{\textwidth}{Y L{5.8cm}}
\toprule
Indicador & Resultado\\
\midrule
Registros de teste & 10\,800\\
Acurácia geral do Random Forest & 98,9907\%\\
Acurácia balanceada & 49,9860\%\\
F1 macro, considerando quatro rótulos & 24,8732\%\\
Acurácia de prever sempre \texttt{low} & 99,0185\%\\
\bottomrule
\end{tabularx}
\source{relatório \texttt{embedded-ai-metrics.json}; referência constante calculada pela proporção da classe \texttt{low}}
\end{table}

\begin{table}[htbp]
\centering
\caption{Representatividade das classes no teste e sensibilidade observada}
\label{tab:classes}
\begin{tabularx}{\textwidth}{L{2.7cm}L{2.4cm}L{3cm}Y}
\toprule
Rótulo & Suporte & Recall & Interpretação\\
\midrule
\texttt{low} & 10\,694 & 99,9719\% & Classe amplamente predominante.\\
\texttt{moderate} & 106 & 0\% & Nenhum exemplo dessa classe foi reconhecido corretamente.\\
\texttt{high} & 0 & Não estimável & Ausente no conjunto de teste.\\
\texttt{critical} & 0 & Não estimável & Ausente no conjunto de teste.\\
\bottomrule
\end{tabularx}
\source{relatório de classificação do modelo compacto}
\end{table}

O conjunto é fortemente desbalanceado. Uma decisão constante pela
classe \texttt{low} alcançaria
$10\,694/10\,800\approx99{,}02\%$, ligeiramente acima da acurácia
do modelo. Portanto, a acurácia geral elevada não demonstra vantagem
preditiva nesse teste. O recall nulo de \texttt{moderate} impede
afirmar que o classificador reconhece adequadamente os níveis
intermediários.

Há também uma diferença de cálculo entre os indicadores. O F1 macro
reportado inclui os quatro rótulos, atribuindo zero a resultados
indefinidos pela convenção \texttt{zero\_division=0}. A acurácia
balanceada corresponde à média dos recalls das classes presentes
no teste. Não se devem interpretar esses percentuais como avaliação
representativa das quatro condições. Para \texttt{high} e
\texttt{critical}, suporte zero significa falta de evidência,
não demonstração de erro ou acerto nessas condições.

\subsection{Verificação de software e da exportação}

Na revisão do relatório, foram executados 87 testes de backend e IA,
sete testes do frontend e o build TypeScript/Vite. As verificações
abrangem validação do contrato, compatibilidade com firmware anterior,
falhas independentes dos sensores, origem do diagnóstico, persistência
CSV e atualização do histórico. A existência de testes aprovados
não constitui validação metrológica ou biológica.

O script de verificação compilou o preditor C++ em ambiente nativo
e o comparou com o modelo Python em 43\,200 registros. Não houve
divergência de classe, e a maior diferença de probabilidade foi de
aproximadamente $3{,}01\times10^{-8}$. Trata-se de evidência de
fidelidade da exportação. Essa comparação inclui dados de treinamento
e, por isso, não é outra estimativa de generalização do modelo.

O registro de compilação da Heltec V2 descrito na
Seção~\ref{sec:modelo} é uma evidência distinta da comparação nativa.
Não foram realizados, nesta revisão documental, gravação de firmware, ensaio de
calibração, medição de consumo ou cronometragem de inferência na placa.
Essas tarefas permanecem no plano experimental.

\subsection{Plano de testes físicos e de integração}

As Tabelas~\ref{tab:testes} e~\ref{tab:testes-integracao} definem procedimentos e critérios a serem
registrados, sem apresentá-los como resultados já alcançados.
Instrumentos de referência devem ter sua identificação e condição
de calibração documentadas. Antes dos ensaios, devem ser registrados
espécie e massa dos grãos, geometria do mini-silo, posição dos
sensores e condições ambientais; esses elementos não estão
caracterizados no material técnico examinado.

\begin{table}[htbp]
\centering\small
\caption{Plano de verificação dos componentes e das falhas}
\label{tab:testes}
\begin{tabularx}{\textwidth}{L{2.5cm}Y Y}
\toprule
Objeto & Procedimento e instrumento & Evidência a registrar\\
\midrule
Alimentação &
Medir entrada USB e barramento de 3,3~V com multímetro; observar quedas
durante transmissão com osciloscópio, se disponível. &
Tensão mínima/máxima, corrente média e picos observáveis; confrontar com
a fonte e o regulador efetivos. Não inferir picos apenas da média.\\
DHT22 &
Comparar temperatura e UR com termo-higrômetro de referência em condições
estáveis e após mudança controlada de ambiente. &
Diferenças entre instrumentos, dispersão, tempo de estabilização e
ocorrência de leituras inválidas. Não confundir concordância pontual com calibração.\\
Sonda e MAX6675 &
Observar repouso, aquecer levemente a ponta com a mão e aguardar
estabilização; comparar com termômetro de referência em pontos adequados
à montagem. Verificar termopar aberto. &
Resposta térmica, diferença em relação à referência e sinalização
\texttt{probe\_ok=false}. Um valor numérico isolado não comprova exatidão.\\
SW-420 &
Aplicar perturbações reproduzíveis na estrutura; observar DO e transições
com analisador lógico, quando disponível. &
Polaridade, limiar ajustado, janela real e repetibilidade. Não converter
contagem em intensidade, RPM ou estado de ventilação.\\
Falhas e IA &
Introduzir falhas independentes de DHT22 e termopar, preservando conexões
seguras; conferir o payload e a serial. &
Medições inválidas null, classe \texttt{unavailable}, score/confiança
null e continuidade da telemetria válida de vibração.\\
\bottomrule
\end{tabularx}
\source{elaboração própria; plano de trabalho, sem resultado físico presumido}
\end{table}

\begin{table}[htbp]
\centering\small
\caption{Plano de verificação da integração e da validade do classificador}
\label{tab:testes-integracao}
\begin{tabularx}{\textwidth}{L{2.5cm}Y Y}
\toprule
Objeto & Procedimento e instrumento & Evidência a registrar\\
\midrule
Execução local &
Com sensores válidos, observar serial e interromper o Wi-Fi;
registrar a duração da inferência e memória livre em execução. &
Classificação local sem depender da API; distribuição de tempos
de \texttt{inference\_us}, reinicializações e comportamento do heap.\\
Rede e API &
Interromper/restabelecer Wi-Fi e servidor; enviar payloads válidos
e inválidos em ambiente de teste isolado. &
Retornos 201/422/503 quando aplicáveis, estados de frescor,
reconexão e inexistência de classificação reapresentada como nova.\\
CSV e interface &
Reiniciar a API, consultar o histórico, exportar o CSV e interromper a
rede do navegador durante uma coleta de teste. &
Persistência das linhas aceitas, indicadores de validade, origem
\texttt{esp32}, preservação de dados antigos e atualização de novos pontos.\\
Validade do classificador &
Definir protocolo de coleta por lotes independentes, com referências
físicas e, quando cabível, análise microbiológica por laboratório habilitado. &
Rótulos externos às previsões, matriz de confusão, recall por classe,
incerteza das métricas e avaliação de lotes não usados no ajuste.\\
\bottomrule
\end{tabularx}
\source{elaboração própria; plano de trabalho, sem resultado físico presumido}
\end{table}

Para investigar fungos, os rótulos futuros não devem ser produzidos
apenas pela própria temperatura ou pela classe prevista pelo modelo:
isso criaria uma avaliação circular. A aquisição ambiental e a
determinação de referência devem ser independentes. O planejamento
também deve considerar variações de tipo de grão, teor de água,
localização da sonda e fontes externas de aquecimento.
\FloatBarrier
